% =====================================================================
%  델파이 설문 문항 조판 매크로
%  (설문지와 논문 부록에서 공통으로 사용 — 패키지 로드 후 \input)
% =====================================================================
\usepackage{longtable,array,colortbl,xcolor}
\usepackage[most]{tcolorbox}

\definecolor{DelphiDomain}{HTML}{1F3A5F}
\definecolor{DelphiModule}{HTML}{E8EEF5}
\definecolor{DelphiNew}{HTML}{B45309}
\definecolor{DelphiNewBg}{HTML}{FFF7ED}
\definecolor{DelphiRule}{HTML}{B8C2CC}

% ---- 설정값 (문서에서 \renewcommand / \setboolean 으로 변경) ----
\newif\ifShowProposedItems  \ShowProposedItemstrue   % [신규] 제안 문항 표시
\newif\ifShowCommentColumn  \ShowCommentColumntrue   % 문항별 의견 칸
\newif\ifShowDomainRating   \ShowDomainRatingtrue    % 영역 전체 타당성 문항
\newif\ifShowOpenBox        \ShowOpenBoxtrue         % 영역별 개방형 의견란
\providecommand{\DelphiRound}{1}                     % 조사 차수 (2 이상이면 이전 차수 결과 열 표시)

% ---- 열 너비 ----
\newlength{\DelphiNoW}      \setlength{\DelphiNoW}{10mm}
\newlength{\DelphiRateW}    \setlength{\DelphiRateW}{7.5mm}
\newlength{\DelphiCommentW} \setlength{\DelphiCommentW}{28mm}
\newlength{\DelphiPrevW}    \setlength{\DelphiPrevW}{17mm}
\newlength{\DelphiItemW}

% ---- 이전 차수 통계 (analysis/delphi_analysis.py 가 생성) ----
%   \SetPrevStat{문항번호}{중앙값}{Q1}{Q3}{CVR}
\newcommand{\SetPrevStat}[5]{%
  \expandafter\gdef\csname delphi@prev@#1\endcsname{%
    \textbf{#2}\newline{\tiny(#3--#4)}\newline{\tiny CVR #5}}}
\newcommand{\PrevStat}[1]{%
  \ifcsname delphi@prev@#1\endcsname\csname delphi@prev@#1\endcsname\else--\fi}

% ---- 내부 구현 ----
\makeatletter
\newcount\delphi@ncols
\newif\ifdelphi@prev
\newcommand{\delphi@rows}{}
\newcommand{\delphi@add}[1]{\g@addto@macro\delphi@rows{#1}}
\newcommand{\delphi@rate}{%
  & \centering ① & \centering ② & \centering ③ & \centering ④ & \centering ⑤}

\newcommand{\delphi@setup}{%
  \ifnum\DelphiRound>1 \delphi@prevtrue\else\delphi@prevfalse\fi
  \delphi@ncols=7
  \setlength{\DelphiItemW}{\dimexpr\linewidth-\DelphiNoW-5\DelphiRateW-14\tabcolsep\relax}%
  \ifShowCommentColumn
    \advance\delphi@ncols 1
    \addtolength{\DelphiItemW}{-\dimexpr\DelphiCommentW+2\tabcolsep\relax}\fi
  \ifdelphi@prev
    \advance\delphi@ncols 1
    \addtolength{\DelphiItemW}{-\dimexpr\DelphiPrevW+2\tabcolsep\relax}\fi
  \delphi@build}

% 열 정의와 머리행을 조건에 따라 미리 조립 (표 안에서 \if 를 쓰지 않기 위함)
\newcommand{\delphi@spec}{}
\newcommand{\delphi@headrow}{}
\newcommand{\delphi@hd}[1]{\color{white}\bfseries #1}
\newcommand{\delphi@build}{%
  \gdef\delphi@spec{>{\centering\arraybackslash\small}p{\DelphiNoW}
                    >{\raggedright\arraybackslash\small}p{\DelphiItemW}}%
  \gdef\delphi@headrow{\rowcolor{DelphiDomain}\delphi@hd{번호} & \delphi@hd{세부 구성요소}}%
  \ifdelphi@prev
    \g@addto@macro\delphi@spec{>{\centering\arraybackslash\footnotesize}p{\DelphiPrevW}}%
    \xdef\delphi@prevlabel{\the\numexpr\DelphiRound-1\relax}%
    \g@addto@macro\delphi@headrow{& \delphi@hd{\scriptsize \delphi@prevlabel 차 결과}\newline
                                    \color{white}\tiny Mdn (Q1--Q3)}%
  \fi
  \g@addto@macro\delphi@spec{*{5}{>{\centering\arraybackslash}p{\DelphiRateW}}}%
  \g@addto@macro\delphi@headrow{%
    & \delphi@hd{\scriptsize 전혀\newline 부적절\newline (1)}
    & \delphi@hd{\scriptsize 부적절\newline ~\newline (2)}
    & \delphi@hd{\scriptsize 보통\newline ~\newline (3)}
    & \delphi@hd{\scriptsize 적절\newline ~\newline (4)}
    & \delphi@hd{\scriptsize 매우\newline 적절\newline (5)}}%
  \ifShowCommentColumn
    \g@addto@macro\delphi@spec{>{\raggedright\arraybackslash\small}p{\DelphiCommentW}}%
    \g@addto@macro\delphi@headrow{& \delphi@hd{수정·보완 의견}}%
  \fi}

\newcommand{\delphi@head}{%
  \arrayrulecolor{DelphiDomain}\hline
  \delphi@headrow \\ \hline
  \arrayrulecolor{DelphiRule}}

% ---- 사용자 매크로 ----
\newcommand{\delphi@domnum}{}
\newcommand{\delphi@domname}{}
\newcommand{\Domain}[3]{%
  \gdef\delphi@domnum{#1}\gdef\delphi@domname{#2}%
  \gdef\delphi@rows{}%
  \delphi@setup
  \DomainHeading{#1}{#2}{#3}}

% 영역 제목 모양 (설문지/부록에서 재정의 가능)
\providecommand{\DomainHeading}[3]{%
  \par\medskip
  \begin{tcolorbox}[colback=DelphiModule,colframe=DelphiDomain,boxrule=0.6pt,
                    arc=1mm,left=2mm,right=2mm,top=1mm,bottom=1mm,
                    title={\bfseries 영역 #1.\ #2},fonttitle=\normalsize]
    \small #3
  \end{tcolorbox}}

\newcommand{\delphi@modulerow}[3]{%
  \delphi@add{\rowcolor{DelphiModule}
    \multicolumn{\the\delphi@ncols}{l}{\small\bfseries #1\ #2#3}\\ \hline}}
\newcommand{\Module}[2]{\delphi@modulerow{#1}{#2}{}}
\newcommand{\NewModule}[2]{%
  \ifShowProposedItems
    \delphi@modulerow{#1}{#2}{\ \textcolor{DelphiNew}{\scriptsize[신규 제안]}}\fi}

\newcommand{\delphi@itemrow}[4]{%
  \delphi@add{#4 #1 & \textbf{#2:} #3}%
  \ifdelphi@prev\delphi@add{& \PrevStat{#1}}\fi
  \delphi@add{\delphi@rate}%
  \ifShowCommentColumn\delphi@add{&}\fi
  \delphi@add{\\ \hline}}
\newcommand{\Item}[3]{\delphi@itemrow{#1}{#2}{#3}{}}
\newcommand{\NewItem}[3]{%
  \ifShowProposedItems
    \delphi@itemrow{#1}{#2 \textcolor{DelphiNew}{\scriptsize[신규]}}{#3}{\rowcolor{DelphiNewBg}}\fi}

\newcommand{\EndDomain}{%
  \ifShowDomainRating
    \delphi@itemrow{D\delphi@domnum}{영역 전체}%
      {위 중분류·세부 구성요소가 이 영역을 적절하게 구성하고 있다}{\rowcolor{DelphiModule}}%
  \fi
  \begingroup
  \setlength{\LTpre}{2pt}\setlength{\LTpost}{4pt}%
  \renewcommand{\arraystretch}{1.25}%
  \expandafter\longtable\expandafter{\delphi@spec}
    \delphi@head \endfirsthead
    \delphi@head \endhead
    \delphi@rows
  \endlongtable
  \endgroup
  \ifShowOpenBox\OpenCommentBox{\delphi@domnum}\fi}

% ---- 영역별 개방형 의견란 ----
\newcommand{\delphi@blank}{\par\vspace{6mm}}
\providecommand{\OpenCommentBox}[1]{%
  \begin{tcolorbox}[colback=white,colframe=DelphiRule,boxrule=0.5pt,arc=1mm,
                    left=2mm,right=2mm,top=1mm,bottom=1mm,breakable,
                    title={\small\bfseries [개방형 의견란] 영역 #1},
                    coltitle=DelphiDomain,colbacktitle=DelphiModule]
    \small
    \textbf{추가}가 필요한 구성요소:\delphi@blank
    \textbf{수정}이 필요한 구성요소 (번호 및 수정안):\delphi@blank
    \textbf{삭제}가 필요한 구성요소 (번호 및 사유):\delphi@blank
  \end{tcolorbox}}
\makeatother
